% ===================================================================
%  CADRO N.º 16 — O gran almacén. — O bazar. — Os xoguetes.
%  Traduction 2026 d'apres texto/io, controlee sur texto/fr, texto/es
%  et texto/pt. Consigne : voir texto/gl/10-cadro-01.tex.
%
%  Le tableau d'astronomie (bloc 15-1) porte des renvois a LETTRES,
%  (a) a (f), graves sur la planche elle-meme : ils se rangent sous
%  le numero 85, comme le dit gravuri/literi.json. Ce bloc est aussi
%  le troisieme ecart declare de renvoji.py.
%
%  LA MENAGERIE DE CARTON DONNE ICI LE MEME RESULTAT QU'EN
%  IRLANDAIS : les betes qu'on a vues gardent leur mot — « raposo »
%  le renard, « lobo » le loup, « rato » la souris, « morcego » la
%  chauve-souris, « donicela » la belette, « curuxa » la chouette,
%  « merlo » le merle — et les betes qu'on n'a jamais vues arrivent
%  avec le
%  livre d'images : « hipopótamo », « dromedario », « leopardo »,
%  « hiena », « crocodilo », « pantera », « avestruz ».
%
%  ET LE GALICIEN AJOUTE UN CAS QUE L'IRLANDAIS N'AVAIT PAS : la
%  belette. « Donicela » n'est ni « comadreja » ni « doninha » ; les
%  trois langues ont trois mots pour une bete que toutes trois ont
%  vue. C'est le dernier temoignage de la colonne, et il vaut d'etre
%  pose a cote du premier : les six mots du corps au tableau 2
%  divergeaient de la meme facon. CE QUI SE VOIT SE NOMME
%  SEPAREMENT ; ce qui se lit se nomme ensemble.
%
%  DERNIER TABLEAU DE LA COLONNE GALICIENNE.
% ===================================================================

%%K t16-apar-1 apar
\VUsaut{4.0mm}
\VUtitre{15.0pt}{60}{CADRO N.º 16}
\VUsaut{2.0mm}
\VUpk{13.2pt}{40}{O gran almacén. --- O bazar.}
\VUpk{13.2pt}{40}{Os xoguetes.}
\VUsaut{3.4mm}

%%K t16-01-1 p
1. --- O ano novo achégase. Os grandes almacéns prepararon os seus escaparates de
\VUgras{xoguetes} \textsuperscript{(1)} e de agasallos de ano novo.

%%K t16-01-2 p
Pedinlle ao meu avó que me levase ao bazar para ver esa bela exposición, e el
consentiu.

%%K t16-02-1 p
2. --- Primeiro paramos diante duns belos panoplios que teñen un \VUgras{fusil,} un
\VUgras{quepis,} un \VUgras{sabre,} un \VUgras{cinguidoiro de espada} e un par de
\VUgras{dragonas,} ou un \VUgras{casco,} unha \VUgras{coiraza}
\textsuperscript{(3)} e unha \VUgras{pistola} \textsuperscript{(4)}. Todo iso
interesoume moito máis ca os \VUgras{proxectores} \textsuperscript{(2)}.

%%K t16-tit-1 sub
\VUsaut{3.0mm}
\VUpk{10.2pt}{40}{Belos espectáculos.}
\VUsaut{2.6mm}

%%K t16-03-1 p
3. --- Na mesma sección había un \VUgras{centinela} \textsuperscript{(5)} na súa
\VUgras{garita}; semellaba montar garda diante dun zoolóxico enteiro de cartón. Cantos
animais había alí: un \VUgras{papagaio} \textsuperscript{(6)} no seu pousadoiro, un
\VUgras{rinoceronte} \textsuperscript{(7)} co corno no nariz, un \VUgras{reno}
\textsuperscript{(8)}, un \VUgras{avestruz} \textsuperscript{(9)}, un \VUgras{mono}
\textsuperscript{(10)} rillando unha noz, un \VUgras{raposo} \textsuperscript{(11)}
que levaba unha galiña, un \VUgras{crocodilo} \textsuperscript{(12)} coas enormes
mandíbulas abertas, un \VUgras{camelo} \textsuperscript{(13)}, un elegante
\VUgras{galgo} \textsuperscript{(14)}, unha \VUgras{pantera}
\textsuperscript{(15)}, e despois dúas bestas que eu non coñecía e que son, segundo
din, unha \VUgras{donicela} \textsuperscript{(16)} e unha \VUgras{hiena}
\textsuperscript{(17)}. Tamén había alí un \VUgras{leopardo} \textsuperscript{(18)} e
un \VUgras{lobo} \textsuperscript{(21)}; moi preto había \VUgras{ratiños}
\textsuperscript{(19)} monísimos e \VUgras{ratos} pequeniños \textsuperscript{(20)}
de caucho. Por último, un \VUgras{hipopótamo} \textsuperscript{(22)} e un
\VUgras{dromedario} \textsuperscript{(23)} xibudo remataban a colección. Quedaría alí
moitas máis horas se non houbese preto un espléndido campamento cheo de
\VUgras{soldados de chumbo} \textsuperscript{(29)} que me chamou a atención.

%%K t16-tit-2 sub
\VUsaut{4.0mm}
\VUpk{10.2pt}{40}{Soldados e guerra.}
\VUsaut{2.8mm}

%%K t16-04-1 p
4. --- O meu avó, que é \VUgras{inválido} \textsuperscript{(24)} e que tivo que deixar
o exército por unha \VUgras{ferida} que lle valeu a \VUgras{medalla militar,} gusta
moito de me contar as \VUgras{campañas} nas que participou. Estaba feliz explicando os
distintos movementos do pequeno exército en miniatura. Un grupo de soldados de verdade
que visitaban o bazar --- un \VUgras{soldado de enxeñaría} \textsuperscript{(25)}, un
\VUgras{cazador alpino} \textsuperscript{(26)} coa cabeza cuberta cunha
\VUgras{boina,} un \VUgras{sarxento maior} \textsuperscript{(27)} de infantaría e un
\VUgras{sarxento de artillaría} \textsuperscript{(28)} cun \VUgras{galón} de ouro na
manga --- parou preto de nós para escoitar as explicacións.

%%K t16-05-1 p
5. --- O meu avó, moi orgulloso de ter un auditorio tan brillante, improvisou un plan
de batalla completo.

%%K t16-05-2 p
A cidade fortificada, coas súas \VUgras{murallas} e os seus \VUgras{fosos,} estaba
\VUgras{asediada} polo \VUgras{exército inimigo}, que, despois dun longo
\VUgras{bombardeo}, estaba listo para o \VUgras{asalto}. Pero os \VUgras{asediados},
forzados pola fame, tentaron unha \VUgras{saída} contra o \VUgras{campamento} dos
\VUgras{asediantes}. Despois de rexeitaren o \VUgras{ataque} cun \VUgras{combate}
teimoso arredor do \VUgras{campamento}, os asediantes \VUgras{vencedores} lanzaron os
seus \VUgras{escuadróns} para perseguiren os atacantes \VUgras{vencidos}. Estes
\VUgras{retiráronse} cubrindo o \VUgras{campo de batalla} de \VUgras{mortos} e de
\VUgras{feridos}. Uns \VUgras{padioleiros} levaron estes últimos á \VUgras{ambulancia}
para que o \VUgras{cirurxián} os atendese.

%%K t16-05-3 p
A pesar da resistencia heroica dalgúns \VUgras{batallóns} que formaran un
\VUgras{cadro}, a \VUgras{derrota} foi completa, o resto do exército \VUgras{fuxiu}
deixando moitos \VUgras{prisioneiros}. O inimigo foi rematar a súa \VUgras{vitoria}
ocupando a cidade. Se cadra ese será o fin da campaña e, despois dun
\VUgras{armisticio}, será posible asinar un \VUgras{tratado de paz}.

%%K t16-tit-3 sub
\VUsaut{4.4mm}
\VUpk{10.2pt}{40}{Agasallos de ano novo.}
\VUsaut{2.8mm}

%%K t16-06-1 p
6. --- Os soldados novos, que rían escoitando o relato pintoresco do veterano,
pedíronlle algunhas explicacións máis. Pola miña parte, dei a volta á tenda para mirar
unha bela \VUgras{bésta} \textsuperscript{(32)} e un \VUgras{arco}
\textsuperscript{(33)} coa súa \VUgras{frecha} \textsuperscript{(31)}. Alí atopei
outra colección de bestas: dous \VUgras{paxaros cantores} \textsuperscript{(34)} ---
un \VUgras{reiseñor} automático e unha \VUgras{papuxa} que cantaba a plena gorxa ---
e unha serie de animais estraños: un \VUgras{morcego} \textsuperscript{(35)}, unha
\VUgras{boa} \textsuperscript{(36)}, unha \VUgras{tartaruga}
\textsuperscript{(37)}, unha \VUgras{foca} \textsuperscript{(38)}, unha
\VUgras{balea} \textsuperscript{(39)} e un \VUgras{tiburón} \textsuperscript{(40)} de
pano.

%%K t16-07-1 p
7. --- Non parei moito para os contemplar, pois albisquei o que soñaba: un
belísimo \VUgras{teatro de monicreques} \textsuperscript{(41)} con espléndidos
\VUgras{decorados} e un encantador \VUgras{pano} vermello. No \VUgras{escenario}, dous
actores semellaban representar un \VUgras{papel} nunha \VUgras{obra} moi interesante,
pois os pequenos \VUgras{espectadores} de cartón instalados nos palcos ou sentados nas
\VUgras{butacas} e no \VUgras{patio} detrás do \VUgras{director de orquestra}
aplaudían con forza.

%%K t16-07-2 p
Por desgraza, ese teatro custaba moi caro.

%%K t16-08-1 p
8. --- Ademais, escoller os xoguetes era difícil, pois nos escaparates amoreábase toda
clase de xoguetes: \VUgras{Arlequín} \textsuperscript{(42)}, \VUgras{Polichinelo}
\textsuperscript{(43)}, \VUgras{Pierrot} \textsuperscript{(44)}, \VUgras{curuxas}
\textsuperscript{(45)} que batían as ás, \VUgras{merlos} \textsuperscript{(46)} que
asubiaban, \VUgras{cans de augas} que ladraban, un \VUgras{fonógrafo}
\textsuperscript{(47)} e \VUgras{xogos de paciencia.} Un deses xoguetes divertiume
moito: representaba un \VUgras{combate naval} \textsuperscript{(48)} no que un
\VUgras{navío} era atacado por \VUgras{piratas} preto dunha terra plantada de
\VUgras{coqueiros}.

%%K t16-noto-1 noto
\VUnotes{143.2mm}{%
\textit{(*) Véxase o cadro N.º 6.}%
}

%%K t16-09-1 p
9. --- Puiden contemplar moi doadamente todas esas marabillas, pois había pouca xente
no almacén, apenas algúns paseantes, un \VUgras{dragón} \textsuperscript{(49)} e un
\VUgras{cobrador} \textsuperscript{(50)} que presentaba unha \VUgras{letra}
na \VUgras{caixa} principal \textsuperscript{(51)}.

%%K t16-10-1 p
10. --- Pregunteille ao meu avó para que servían os libros postos en estantes detrás da
\VUgras{caixeira}; respondeume que son \VUgras{libros de contabilidade}.

%%K t16-tit-4 sub
\VUsaut{3.6mm}
\VUpk{10.2pt}{40}{Papar moscas arredor da tenda.}
\VUsaut{2.8mm}

%%K t16-11-1 p
11. --- Ao deixar a sección dos xoguetes, fomos á talabartaría botar unha ollada ás
\VUgras{selas} \textsuperscript{(53)}, aos \VUgras{estribos} \textsuperscript{(54)},
ás \VUgras{bridas} \textsuperscript{(55)}, aos \VUgras{bocados}
\textsuperscript{(56)} e ás \VUgras{fustas.} Mentres o \VUgras{xefe de sección}
\textsuperscript{(57)} lle daba algúns informes ao meu avó, eu fun mirar o escaparate
das \VUgras{porcelanas} e dos \VUgras{cristais} \textsuperscript{(58)}, onde hai belas
\VUgras{ensaladeiras} e delicadas \VUgras{teteiras} \textsuperscript{(59)}.

%%K t16-12-1 p
12. --- Para ir ao primeiro andar pasamos detrás do escaparate dos \VUgras{corsés}
\textsuperscript{(60)}, ao lado da tenda de medias, onde estaban expostos uns
\VUgras{calzóns} moi belos \textsuperscript{(63)}. Preto de alí, unha
\VUgras{dependenta} \textsuperscript{(61)} facía escoller a unha \VUgras{compradora}
\textsuperscript{(62)} uns belos \VUgras{tecidos} de seda \textsuperscript{(64)}.

%%K t16-12-2 p
Subindo pola escaleira, reparei en que o almacén está adornado con \VUgras{farois}
xaponeses \textsuperscript{(65)}, \VUgras{parasoles} \textsuperscript{(66)} e enormes
\VUgras{palmeiras} \textsuperscript{(70)}.

%%K t16-13-1 p
13 . --- No primeiro andar non había moito que ver; só mobles (*), \VUgras{caixas
fortes} \textsuperscript{(67)}, \VUgras{cómodas} \textsuperscript{(68)} e
\VUgras{carriños de neno} \textsuperscript{(69)}. Pero a vista de conxunto do almacén
era moi \VUgras{divertida.}

%%K t16-14-1 p
14. --- Aos nosos pés vin os instrumentos de música: \VUgras{arpas}
\textsuperscript{(71)}, \VUgras{guitarras} \textsuperscript{(72)},
\VUgras{bandolinas} \textsuperscript{(73)}, \VUgras{cornetíns de pistóns}
\textsuperscript{(74)}, \VUgras{clarinetes} \textsuperscript{(75)}, violíns co seu
\VUgras{arco} \textsuperscript{(76)} e mesmo un \VUgras{bombo}
\textsuperscript{(77)}. Un pouco máis lonxe, na sección do papel, un
\VUgras{estudante} \textsuperscript{(78)} regateaba a un \VUgras{dependente}
\textsuperscript{(79)} unha carpeta de cadernos. Preto deles distinguín un belo
\VUgras{abecedario} \textsuperscript{(80)}.

%%K t16-14-2 p
No medio do almacén erguérase unha alta \VUgras{árbore de Nadal}
\textsuperscript{(81)} toda provista de candeas, de miudezas e de toda clase de
xoguetes: \VUgras{cabalos de madeira} \textsuperscript{(82)}, \VUgras{bonecas}
\textsuperscript{(83)}, etc.

%%K t16-14-3 p
A \VUgras{perfumaría} \textsuperscript{(84)}, cos seus \VUgras{dentífricos} e os seus
\VUgras{perfumes} variados, espallaba un moi bo recendo que chegaba ata nós.

%%K t16-tit-5 sub
\VUsaut{3.4mm}
\VUpk{10.2pt}{40}{Mercamos algo.}
\VUsaut{2.8mm}

%%K t16-15-1 p
15. --- Como o meu avó comezaba a xulgar longo de máis o tempo, baixamos pola escaleira
grande. Parou un intre diante dun \VUgras{cadro de astronomía} \textsuperscript{(85)}
que representaba, díxome, \VUgras{cometas} \textsuperscript{(a)}, \VUgras{eclipses de
lúa e de sol} \textsuperscript{(b)}, unha rosa dos ventos cos \VUgras{catro puntos
cardinais} \textsuperscript{(c)} \VUgras{(Norte, Sur, Leste e Oeste)}, un mapa dos
dous \VUgras{hemisferios} \textsuperscript{(d)}, os \VUgras{planetas}
\textsuperscript{(e)} e o \VUgras{sistema solar} \textsuperscript{(f)}. Eu non entendín
nada de todo iso, e interesoume moito máis a librea con galóns dun \VUgras{mozo de
reparto} \textsuperscript{(86)} que pasaba, e mais os belos \VUgras{abanicos}
\textsuperscript{(87)} que estaban preto de min, ao lado dos \VUgras{moedeiros}
\textsuperscript{(88)} e das \VUgras{carteiras} \textsuperscript{(89)}.

%%K t16-16-1 p
16. --- Admirei tamén no escaparate \VUgras{plumas} \textsuperscript{(90)} e
\VUgras{fibelas de cinto} \textsuperscript{(91)}, e as belas \VUgras{flores
artificiais} \textsuperscript{(94)} dispostas con graza en cestos. As
\VUgras{violetas}, os \VUgras{lirios}, os \VUgras{xacintos}
\textsuperscript{(95)}, os \VUgras{xeranios} \textsuperscript{(96)}, os
\VUgras{lirios brancos} \textsuperscript{(97)} e as \VUgras{capuchinas}
\textsuperscript{(98)} estaban moi ben copiados.

%%K t16-tit-6 sub
\VUsaut{4.4mm}
\VUpk{10.2pt}{40}{Un agasallo do meu avó.}
\VUsaut{2.8mm}

%%K t16-17-1 p
17. --- Pedinlle ao meu avó que me mercase un dos \VUgras{cepillos de dentes}
\textsuperscript{(92)} que estaban nunha caixa preto duns \VUgras{gafitos de pelo}
\textsuperscript{(93)}. Consentiu e fun eu mesmo pagar a miña compra á caixa, preto da
cal se preparaba unha importante \VUgras{expedición} \textsuperscript{(100)} para un
\VUgras{cliente} \textsuperscript{(99)}.

%%K t16-18-1 p
18. --- De súpeto houbo un lixeiro rebumbio entre as persoas paradas preto dos
\VUgras{bronces} artísticos \textsuperscript{(101)}. Un \VUgras{ladrón}
\textsuperscript{(102)} acababa de ser collido no seu delito por un \VUgras{inspector}
\textsuperscript{(103)}, mentres tentaba roubar alguén. Coidouse de facer vir un
policía para o \VUgras{deter,} e xusto cando estabamos a saír levaban o delincuente ao
cárcere.
